% F087｜牛顿切线法快速求近似根｜03_proof.tex
% 继承上传模板的 LaTeX 环境；当前条目的内容独立编写。
\clearpage
\begin{proofbox}
\textbf{\textcolor{CProof}{证明一（门槛最低）：从切线方程推导}}
\par\medskip
\begin{center}
% F087｜tangent 原创简明教学示意图
\begin{tikzpicture}[>=stealth,scale=.82]
\draw[->] (0,0)--(7.6,0) node[right]{$x$};\draw[->] (.8,-1.4)--(.8,3.2) node[above]{$y$};
% 横向比例尺 x_point=8(x-2.65)，纵向比例尺 y_point=0.75(x^2-10)。
\draw[teal,thick,domain=2.76:3.48,samples=45] plot ({8*(\x-2.65)},{.75*(\x*\x-10)});
% 在真实点 x_0=3 的切线；交轴位置 x_1=19/6。
\draw[orange,thick] (2,-1.2)--(6,1.05);
\fill[orange] (2.8,-.75) circle (2pt) node[below left=5pt]{$x_0=3$};
\fill[orange] (4.1333,0) circle(2pt) node[below=7pt]{$x_1\approx3.1667$};
\node at (5.5,2.7){\small 用切线寻找新的交点};
\node[teal] at (5.7,1.9){\small $y=x^2-10$};
\end{tikzpicture}

\end{center}
\medskip
切线方程为 $y-f(x_n)=f'(x_n)(x-x_n)$。令 $y=0$，则 $-f(x_n)=f'(x_n)(x_{n+1}-x_n)$，所以\[\boxed{x_{n+1}=x_n-\frac{f(x_n)}{f'(x_n)}}.\]这里必须要求 $f'(x_n)\ne0$。
\end{proofbox}
\clearpage
\begin{proofbox}
\textbf{\textcolor{CProof}{证明二（最快）：平方根的代数推导}}
\par\medskip
\begin{center}
% F087｜square 原创简明教学示意图
\begin{tikzpicture}[>=stealth,scale=.9]
\node[draw,fill=teal!8,rounded corners,minimum width=7.3cm,minimum height=1.5cm,align=center] at (4,0) {$x_{n+1}=\dfrac12\left(x_n+\dfrac{A}{x_n}\right)$};
\node at (4,-1.2) {\small 正数 $x_n$ 与 $A/x_n$ 的平均数};
\end{tikzpicture}

\end{center}
\medskip
设 $A>0,x_n>0$，将 $f(x)=x^2-A$、$f'(x)=2x$ 代入切线式，整理即得\[x_{n+1}=\frac{x_n^2+A}{2x_n}=\frac12\left(x_n+\frac A{x_n}\right).\]若 $r=\sqrt A$，还可验证 $x_{n+1}-r=\dfrac{(x_n-r)^2}{2x_n}\ge0$。
\end{proofbox}
\clearpage
\begin{proofbox}
\textbf{\textcolor{CProof}{证明三（迁移到立方根）：直接换函数}}
\par\medskip
\begin{center}
% F087｜cube 原创简明教学示意图
\begin{tikzpicture}[>=stealth,scale=.88]
\node[draw,rounded corners,fill=teal!8,minimum width=8cm,minimum height=1.3cm] (a) at (4,0) {$x_{n+1}=\dfrac{2x_n+A/x_n^2}{3}$};
\node at (4,-1.0) {\small 三次根号也能用切线修正，不必背另一套思路};
\end{tikzpicture}

\end{center}
\medskip
设 $A>0,x_n>0$。令 $f(x)=x^3-A$，则 $f'(x)=3x^2$，因此\[x_{n+1}=x_n-\frac{x_n^3-A}{3x_n^2}=\frac{2x_n+A/x_n^2}{3}.\]方法并未变化：\textbf{先写方程，再求切线，最后求横轴交点。}
\end{proofbox}
